\documentclass[10pt,a4paper]{amsart}
\usepackage[margin=19mm]{geometry}
\usepackage{amsmath,amssymb,mathtools}
\usepackage[scr=rsfs,cal=euler]{mathalfa}
\usepackage{xfrac}
\usepackage[unicode,hidelinks,bookmarks=false]{hyperref}
\newcommand{\ie}{i.e.\ }
\newcommand{\cf}{cf.\ }
\newcommand{\resp}{respectively\ }
\newcommand{\Subsection}{\subsection}
\providecommand{\nobreakdash}{\nobreak\hbox{-}}
\setlength{\emergencystretch}{2em}
\multlinegap0pt
\usepackage{amssymb}
\usepackage{float}
\newtheorem{theorem}{Theorem}
\title{Two-phase source and reaction coefficient Stefan type problems}
\author{Targyn A. Nauryz}
\date{Source: arXiv:2607.21388v2 (CC BY 4.0)}
\begin{document}
\maketitle

\noindent\emph{Source and licence.} Two-phase source and reaction coefficient Stefan type problems. Version arXiv:2607.21388v2 (CC BY 4.0), distributed by arXiv under the Creative Commons Attribution 4.0 International licence, http://creativecommons.org/licenses/by/4.0/, as recorded in the arXiv licence metadata for this exact version and shown on the abstract page https://arxiv.org/abs/2607.21388v2. Full licence notice: \texttt{licenses/arxiv-2607.21388-CC-BY-4.0.txt}.\par\smallskip

\noindent\textit{Editorial scope.} Counted unit: the article's theorem thm4 (source0.tex lines 681--689). The article's complete original proof of it is delivered with it (source0.tex lines 690--787, one \texttt{proof} environment of 98 lines). No mathematics is written, repaired or re-derived: the statement and the proof are the article's own lines, in the article's own notation, and no retained line is substituted or re-flowed.  Retained uncounted as the source-local context the proof needs: lines 216--226; lines 238--243.  Nothing was omitted from any retained span, so the package carries no omission record.  No retained line contains a citation, so the wrapper carries no bibliography and no bibliography entry.  No retained line contains a figure environment, an image-inclusion command or drawing code, so no image is delivered and the package declares no asset dependency.  Editorial formatting only, in addition: an AMS \texttt{amsart} preamble that restates the article's own declarations and macros used by the retained lines, namely the environments \texttt{theorem} on separate counters, together with the standard packages those lines need.  The retained environments print in this document as Theorem 1 in order of appearance, whereas the article prints the same items with the numbering of its own full document.  The complete native source of the article as distributed in the arXiv e-print for this exact version is bundled unchanged as \texttt{sources/205872/source0.tex}.
\begin{equation}\label{eq10}
    \begin{cases}
        \dfrac{\partial V_1}{\partial t}-b_1(\xi_1,t)\dfrac{\partial V_1}{\partial \xi_1}-a_1(t)\dfrac{\partial^2 V_1}{\partial\xi_1^2}=R_1(t)F_1(\xi_1,t),& (\xi_1,t)\in\Omega,\\
        \dfrac{\partial V_2}{\partial t}-b_2(\xi_2,t)\dfrac{\partial V_2}{\partial \xi_2}-a_2(t)\dfrac{\partial^2 V_2}{\partial\xi_2^2}=R_2(t)F_2(\xi_2,t),& (\xi_2,t)\in\Omega,\\
        V_1(\xi_1,0)=\Phi_1(\xi_1),&\forall\xi_1\in[0,1],\\
        V_2(\xi_2,0)=\Phi_2(\xi_2),&\forall\xi_2\in[0,1],\\
        V_1(0,t)=V_1(1,t)=0,&\forall t\in[0,T],\\
        V_2(0,t)=V_2(1,t)=0,&\forall t\in[0,T],\\
        -\dfrac{k_1}{s(t)}\dfrac{\partial V_1}{\partial \xi_1}\bigg|_{\xi_1=1}=-\dfrac{k_2}{d-s(t)}\dfrac{\partial V_2}{\partial\xi_2}\bigg|_{\xi_2=0}+\chi(t)+Ls'(t),& \forall t\in[0,T].
    \end{cases}
\end{equation}

\begin{equation}\label{eq11}
    \begin{cases}
        -\phi_i''(\xi_i)=\lambda_i\phi_i(\xi_i),&\xi_i\in(0,1),\\
        \phi_i(0)=\phi_i(1)=0,
    \end{cases}
\end{equation}

\begin{theorem}\label{thm4}
Let suppose $s_m$, $s_{m^{\prime}}$, $s_M$ and $s_{M^{\prime}}$ are positive real numbers such that $s(t)\in C^2(0,T)$, $s^{\prime}(t)\in L^{\infty}(0,1)$ with $0<s_m<s(t)<s_M$, $s_{m^{\prime}}<s^{\prime}(t)<s_{M^{\prime}}$ and $s^{\prime}(t)\in L^{\infty}(0,T)$ for all $t\in[0,T]$. Moreover, initial datas $\Phi_{i0}\in \mathbb{H}_0^1(0,1)$ and $F_i(\xi_i,t)\in L^1(0,T,L^2(\Omega))\cap L^2(0,T,L^2(0,1))$, for the uniquely reconstructed coefficients $R_i(t)$ obtained in Theorem 1, 2, and Theorem 3, problem \eqref{eq10} possesses a unique strong solution pair $(V_1,V_2)$ of the problem \eqref{eq10}, that is,
$$\frac{\partial V_i}{\partial t}-a_i(t)\frac{\partial^2 V_i}{\partial \xi^2}-b(\xi_i,t)\frac{\partial V_i}{\partial \xi}=R_i(t)F_i(\xi_i,t),\quad\text{where}\quad i=1,2\quad\text{in}\quad L^2(0,T;L^2(0,1)),$$ which satisfy the following statements for both $i=1,2$:
\begin{itemize}
    \item[a)] $V_i\in L^2(0,T;\mathbb{H}_0^1(0,1)\cap \mathbb{H}^2(0,1))\cap L^{\infty}(0,T;\mathbb{H}_0^1(0,1))$,
    \item[b)] $\partial_t V_i\in L^2(0,T;L^2(0,1))$,
    \item[c)] $R_i(t)\in C[0,T]$, $R_i(t)>0$ and $0<m_i<R_i(t)< M_i$ with $m_i,M_i>0$ for all $t\geq 0$. 
\end{itemize}
\end{theorem}

\begin{proof}
    \textbf{Existence.} Let $T>0$ and define the subspace $X_m$ spanned by $\{\phi_1,\phi_2,...,\phi_m\}$, where $\phi_j$ with $j=1,2,3,...m$ are assumed to be eigenvectors of the space $\mathbb{H}_0^1(0,1)$ that is the solution of the spectral problem \eqref{eq11}. Thus, we can deduce that if $V_{i,m}\in X_m$ where $i=1,2$, then it can be written as linear combination of the eigenvectors $\phi_m$ belonging to $X_m$ such that
    $$V_{i,m}(\xi_i,t)=\sum\limits_{j=1}^m V_{i,j,m}(t)\phi_j(\xi_i),\quad\quad i=1,2,$$
    where $V_{i,j,m}$ is the solution of the system of ordinary differential equation with boundary condition as follows
    \begin{equation}\label{eq62}
        \begin{cases}
            \left(V_{i,m}^{\prime},\phi\right)-a_i(t)\left(\frac{\partial^2 V_{i,m}}{\partial \xi_i^2},\phi\right)-\left(b(\xi_i,t)\frac{\partial V_{i,m}}{\partial \xi_i},\phi\right)=R_i(t)\left(F_i(\xi_i,t),\phi\right),& \forall\phi\in X_m,\\
            V_{i,m}(0)=\Phi_{i,m}\rightarrow \Phi_i,& \text{in}\quad \mathbb{H}_0^1(\Omega) 
        \end{cases}
    \end{equation}
    where $V_{i,m}^{\prime}=\frac{\partial V_{i,m}}{\partial t}$. The system \eqref{eq62}  has local solution in the interval $(0,T_m)$, by Caratheodory Theorem. To extend the local solution to the interval $(0,T)$ independent of $m$ the
    following a priori estimate is needed:

    \begin{itemize}
        \item [Case 1.] If $\phi=V_{i,m}$ where $i=1,2$, then equation in the system \eqref{eq62} takes the form
        $$\left(V_{i,m}^{\prime},V_{i,m}\right)-a_i(t)\left(\frac{\partial^2 V_{i,m}}{\partial \xi_i^2},V_{i,m}\right)-\left(b(\xi_i,t)\frac{\partial V_{i,m}}{\partial \xi_i},V_{i,m}\right)=R_i(t)\left(F_i(\xi_i,t),V_{i,m}\right).$$
        The first term gives us
        $$\left(V_{i,m}^{\prime},V_{i,m}\right)=\int_0^1\frac{\partial V_{i,m}}{\partial t}V_{i,m}d\xi_i=\int_0^1\frac{d}{dt}\left(\frac{1}{2}V_{i,m}\right)^2d\xi_i=\frac{1}{2}\frac{d}{dt}\int_0^1 V_{i,m}^2d\xi_i=\frac{1}{2}\frac{d}{dt}\left|V_{i,m}\right|_{L^2(0,1)}^2.$$
        By integration by parts and boundary conditions in \eqref{eq10}, we estimate the following
        $$a_i(t)\left(\frac{\partial^2 V_{i,m}}{\partial \xi_i^2},V_{i,m}\right)=-a_i(t)\int_0^1 \left(\frac{\partial V_{i,m}}{\partial \xi_i}\right)^2 d\xi_i=-a_i(t)||V_{i,m}||^2_{\mathbb{H}_0^1(0,1)}$$
        and analogously for the third term we get
        $$\left(b(\xi_i,t)\frac{\partial V_{i,m}}{\partial \xi_i},V_{i,m}\right)=\int_0^1 b(\xi_i,t)\frac{\partial V_{i,m}}{\partial \xi_i}V_{i,m}d\xi_i=-\frac{s^{\prime}(t)}{2s(t)}\int_0^1 V_{i,m}^2 d\xi_i=-\frac{s^{\prime}(t)}{2s(t)}\left|V_{i,m}\right|^2_{L^2(0,1)}.$$
        Applying these estimations and Holder inequality, we can provide the following inequality
        $$\frac{1}{2}\frac{d}{dt}|V_{i,m}|_{L^2(0,1)}^2+a_i(t)||V_{i,m}||_{\mathbb{H}_0^1(0,1)}^2+\frac{s^{\prime}(t)}{2s(t)}\left|V_{i,m}\right|^2_{L^2(0,1)}\leq \frac{1}{2}R_i(t)\left(|F_i|_{L^2(0,1)}^2+|V_{i,m}|_{L^2(0,1)}^2\right).$$
        Integrating over $[0,t)$ with $t\in[0,T_m)$ and assuming $0<s_m<s(t)<s_M$, $s_{m^{\prime}}<s^{\prime}(t)<s_{M^{\prime}}$, $|a_i(t)|<N_i$  and (c) hold, then again by integration by parts we have
        $$|V_{i,m}|_{L^2(0,1)}^2+\int_0^t ||V_{i,m}||_{\mathbb{H}_0^1(0,1)}^2d\tau\leq\frac{1}{2c_1}|V_{i,m}(0)|_{L^2(0,1)}^2+\frac{c_2}{c_1}|F_{i}|_{L^2(0,1)}2+\frac{c_3}{c_1}\int_0^T |V_{i,m}|_{L^2(0,1)}^2d\tau,$$
        where $c_1=\min\{\frac{1}{2}, N_i\}$, $c_2=\frac{M_i}{2}$ and $c_3=\frac{M_i}{2}+\frac{s_{M^{\prime}}}{2s_m}$. Then according to the Gronwall's inequality we can deduce that
        $$V_{i,m}\;\text{ is bounded in }\;L^{\infty}(0,T; L^2(0,1))\quad\text{and}\quad V_{i,m}\;\text{ is bounded in }\; L^2(0,T;\mathbb{H}_0^1(0,1)).$$

        \item[Case 2.] If we take $\phi=V_{i,m}^{\prime}=\frac{d}{dt}V_{i,m}$ where $i=1,2$, then the system \eqref{eq62} becomes
        $$\left(V_{i,m}^{\prime},V_{i,m}^{\prime}\right)-a_i(t)\left(\frac{\partial^2 V_{i,m}}{\partial \xi_i^2},V_{i,m}^{\prime}\right)-\left(b(\xi_i,t)\frac{\partial V_{i,m}}{\partial \xi_i},V_{i,m}^{\prime}\right)=R_i(t)\left(F_i(\xi_i,t),V_{i,m}^{\prime}\right),$$
        for all $V_{i,m}^{\prime}\in X_m$.
        Therefore, for the all terms on the left-hand side of this identity we have the following verifications
        $$\left(V_{i,m}^{\prime},V_{i,m}^{\prime}\right)=\int_0^1\left(\frac{dV_{i,m}}{dt}\right)^2 d\xi_i=|V_{i,m}^{\prime}|_{L^2(0,1)}^2,$$
        $$a_i(t)\left(\frac{\partial^2 V_{i,m}}{\partial \xi_i^2},V_{i,m}^{\prime}\right)=a_i(t)\int_0^1\frac{\partial^2V_{i,m}}{\partial\xi_i^2}\frac{dV_{i,m}}{dt}d\xi_i=a_i(t)\frac{d}{dt}\int_0^1\frac{\partial^2V_{i,m}}{\partial\xi_i^2} V_{i,m}d\xi_i$$
        $$=-\frac{a_i(t)}{2}\frac{d}{dt}\int_0^1\left(\frac{\partial V_{i,m}}{\partial\xi_i}\right)^2d\xi_i=-\frac{a_i(t)}{2}\frac{d}{dt}||V_{i,m}||_{\mathbb{H}_0^1(0,1)}^2,$$
        and 
        $$\left(b(\xi_i,t)\frac{\partial V_{i,m}}{\partial \xi_i},V_{i,m}^{\prime}\right)=\int_0^1 b(\xi_i,t)\frac{\partial V_{i,m}}{\partial \xi_i}V_{i,m}^{\prime}d\xi_i=-\frac{s^{\prime}(t)}{s(t)}\int_0^1 V_{i,m}^{\prime} V_{i,m}d\xi_i$$
        $$-\int_0^1 b(\xi_i,t)\frac{\partial V_{i,m^{\prime}}}{\partial \xi_i}V_{i,m}d\xi_i=-\frac{s^{\prime}(t)}{2s(t)}\frac{d}{dt}\int_0^1 V_{i,m}^2d\xi_i=-\frac{s^{\prime}(t)}{2s(t)}\frac{d}{dt}|V_{i,m}|_{L^2(0,1)}^2,$$
        analogously as previous case 1, respect to Holder inequality we obtain
        $$\left(1-\frac{1}{2}R_i(t)\right)|V_{i,m}^{\prime}|_{L^2(0,1)}^2+\frac{a_i(t)}{2}\frac{d}{dt}||V_{i,m}||_{\mathbb{H}_0^1(0,1)}^2+\frac{s^{\prime}(t)}{2s(t)}\frac{d}{dt}|V_{i,m}|_{L^2(0,1)}^2\leq\frac{R_i(t)}{2}|F_i|_{L^2(0,1)}^2.$$
        Hence, employing integrating on $[0,t)$ with $t\in[0,T_m)$ and according to integration by part we have
        $$\int_0^t |V_{i,m}|_{L^2(0,1)}^2d\tau+||V_{i,m}(t)||_{\mathbb{H}_0^1(0,1)}^2+|V_{i,m}(t)|_{L^2(0,1)}^2$$
        $$\leq\frac{c_5}{c_4}||V_{i,m}(0)||_{\mathbb{H}_0^1(0,1)}^2+\frac{c_6}{c_4}|V_{i,m}(0)|_{L^2(0,1)}^2+\frac{c_2}{c_4}|F_i|_{L^2(0,1)}^2,$$
        where $c_5=\frac{N_i}{2}$, $c_6=\frac{s_{M^{\prime}}}{2s_m}T$ and $c_4=\min\{1-\frac{m_i}{2},\frac{N_i}{2}, \frac{s_{M^{\prime}}}{2s_m}\}$. Thus, Gronwall's Lemma guarantees that
        $$V_{i,m}\;\text{ is bounded in }\; L^{\infty}(0,T,\mathbb{H}_0^1(0,1)),\quad\quad V_{i,m}^{\prime}\;\text{ is bounded in }\; L^2(0,T,L^2(0,1)).$$
        \item[Case 3.] Assume that $\phi=-\frac{\partial^2 V_{i,m}}{\partial\xi_i^2}$, then we can rewrite equation in the system \eqref{eq62} as follows
        $$-\left(V_{i,m}^{\prime},\frac{\partial^2 V_{i,m}}{\partial\xi_i^2}\right)+a_i(t)\left(\frac{\partial^2 V_{i,m}}{\partial \xi_i^2},\frac{\partial^2 V_{i,m}}{\partial\xi_i^2}\right)+\left(b(\xi_i,t)\frac{\partial V_{i,m}}{\partial \xi_i},\frac{\partial^2 V_{i,m}}{\partial\xi_i^2}\right)$$
        $$=-R_i(t)\left(F_i(\xi_i,t),\frac{\partial^2 V_{i,m}}{\partial\xi_i^2}\right).$$
        For each term of the left-hand side by integration by parts we can provide the following evaluations
        $$\left(V_{i,m}^{\prime},\frac{\partial^2 V_{i,m}}{\partial\xi_i^2}\right)=\int_0^1 V_{i,m}^{\prime}\frac{\partial^2 V_{i,m}}{\partial\xi_i^2}d\xi_i=-\int_0^1\frac{\partial V_{i,m}^{\prime}}{\partial\xi_i}\frac{\partial V_{i,m}}{\partial\xi_i}d\xi_i=-\frac{1}{2}\frac{d}{dt}\int_0^1\left(\frac{\partial V_{i,m}}{\partial\xi_i}\right)^2d\xi_i$$
        $$=-\frac{1}{2}\frac{d}{dt}||V_{i,m}||_{\mathbb{H}_0^1(0,1)}^2,$$
        $$\left(\frac{\partial^2 V_{i,m}}{\partial \xi_i^2},\frac{\partial^2 V_{i,m}}{\partial\xi_i^2}\right)=\int_0^1\left(\frac{\partial^2V_{i,m}}{\partial\xi_i^2}\right)^2d\xi_i=\left|\frac{\partial^2V_{i,m}}{\partial\xi_i^2}\right|_{L^2(0,1)}^2,$$
        $$\left(b(\xi_i,t)\frac{\partial V_{i,m}}{\partial \xi_i},\frac{\partial^2 V_{i,m}}{\partial\xi_i^2}\right)=\int_0^1 b(\xi_i,t)\frac{\partial V_{i,m}}{\partial\xi_i}\frac{\partial^2 V_{i,m}}{\partial\xi_i^2}d\xi_i=-\frac{s^{\prime}(t)}{2s(t)}\int_0^1\left(\frac{\partial V_{i,m}}{\partial\xi_i}\right)^2d\xi_i$$
        $$=-\frac{s^{\prime}(t)}{2s(t)}||V_{i,m}||_{\mathbb{H}_0^1(0,1)}^2.$$
        Then, by Holder inequality we have
        $$\frac{1}{2}\frac{d}{dt}||V_{i,m}||_{\mathbb{H}_0^1(0,1)}^2+a_i(t)\left|\frac{\partial^2V_{i,m}}{\partial\xi_i^2}\right|_{L^2(0,1)}^2-\frac{s^{\prime}(t)}{2s(t)}||V_{i,m}||_{\mathbb{H}_0^1(0,1)}^2$$
        $$\leq\frac{1}{2}R_i(t)|F_i|_{L^2(0,1)}^2+\frac{1}{2}R_i(t)\left|\frac{\partial^2V_{i,m}}{\partial\xi_i^2}\right|_{L^2(0,1)}^2.$$
        Thus, similarly as previous cases assuming all conditions hold, we establish the following result
        $$||V_{i,m}(t)||_{\mathbb{H}_0^1(0,1)}^2+\int_0^t\left|\frac{\partial^2V_{i,m}}{\partial\xi_i^2}\right|d\tau\leq \frac{c_2}{c_7}|F_i|_{L^2(0,1)}^2+\frac{c_6}{c_7}\int_0^T||V_{i,m}||_{\mathbb{H}_0^1(0,1)}^2d\tau,$$
        where $c_2$, $c_6$ are mentioned in previous cases and $c_7=\min\{\frac{1}{2},N_i-\frac{M_i}{2}\}$. Hence, respect to Gronwall's inequality we conclude that
        $$V_{i,m}\;\text{ is bounded in }\; L^{\infty}(0,T,\mathbb{H}_0^1(0,1)),\quad\quad \frac{\partial^2V_{i,m}}{\partial\xi_i^2}\;\text{ is bounded in }\; L^2(0,T,L^2(0,1)).$$
    \end{itemize}
    \textbf{Uniqueness.} Suppose that the problem admits two global strong solution pairs
    $$(V_{i,m}^{(1)},R_i^{(1)})\quad \text{and}\quad (V_{i,m}^{(2)},R_i^{(2)})\quad i=1,2$$
    corresponding to the same initial, boundary, and overdetermination data. Let us define
    $$W_{i,m}(\xi_i,t)=V_{i,m}^{(1)}(\xi_i,t)-V_{i,m}^{(2)}(\xi_i,t)\quad\text{and}\quad \rho_i(t)=R_i^{(1)}(t)-R_i^{(2)}(t).$$
    Subtracting the two nonhomogeneous equations gives
    \begin{equation}\label{eq63}
        \frac{\partial W_{i,m}}{\partial t}-b(\xi_i,t)\frac{\partial W_{i,m}}{\partial \xi_i}-a_i(t)\frac{\partial^2 W_{i,m}}{\partial\xi_i^2}=\rho_i(t)F_i(\xi_i,t)
    \end{equation}
    and assume that heat source coefficient from overdetermination conditions can be determined in the Volterra integral form
    $$R_i(t)=G_i(t)+\int_0^tK_i(t,\tau)R_i(\tau)d\tau,\quad i=1,2$$ 
    where $G_i(t)$ and $K_i(t,\tau)$ are continuous and bounded functions in $[0,T]$, substracting $R_i^{(1)}$ and $R_i^{(2)}$, it which implies
    \begin{equation}\label{eq64}
        \rho_i(t)=\int_0^t K_i(t,\tau)\rho_i(\tau)d\tau.
    \end{equation}
    If we assume that $|K_i(t,\tau)|<C_i$ holds, then it follows that $|\rho_i(t)|\leq C_i\int_0^t |\rho_i(\tau)|d\tau.$
    By Gronwall's inequality, we obtain that $\rho_i(t)\equiv 0$ for $t\in[0,T]$. Therefore, $R_i^{(1)}(t)=R_i^{(2)}(t)$.
    
    Thus, the differenial equation is also nonhomogeneous unless $\rho_i=0$. Multiplying both sides \eqref{eq63} by $W_{i,m}$ and integrate over $(0,1)$:
    $$\int_0^1\frac{\partial W_{i,m}}{\partial t}W_{i,m}d\xi_i-\int_0^1b(\xi_i,t)\frac{\partial W_{i,m}}{\partial \xi_i} W_{i,m}d\xi_i-a_i(t)\int_0^1\frac{\partial^2 W_{i,m}}{\partial\xi_i^2}W_{i,m}d\xi_i$$
    $$=\rho_i(t)\int_0^1F_i(\xi_i,t)W_{i,m}d\xi_i.$$
    Again, for the each term of the left-hand side we have
    $$\int_0^1\frac{\partial W_{i,m}}{\partial t}W_{i,m}d\xi_i=\frac{1}{2}\frac{d}{dt}||W_{i,m}(t)||_{L^2(0,1)}^2,\quad -a_i(t)\int_0^1\frac{\partial^2 W_{i,m}}{\partial\xi_i^2}W_{i,m}d\xi_i=a_i(t)\left\|\frac{\partial W_{i,m}}{\partial\xi_i}(t)\right\|_{L^2(0,1)}^2,$$
    $$\int_0^1b(\xi_i,t)\frac{\partial W_{i,m}}{\partial \xi_i} W_{i,m}d\xi_i=\frac{1}{2}\int_0^1\frac{\partial b}{\partial \xi_i}W_{i,m}^2d\xi_i.$$
    The right-hand side is estimated by Cauchy–Schwarz and Young inequalities:
    $$|\rho_i(t)(F_i,W_{i,m})_{L^2(0,1)}|\leq |\rho_i(t)|||F_i(t)||_{L^{2}(0,1)}||W_{i,m}(t)||_{L^2(0,1)}$$
    $$\leq \frac{1}{2}|\rho_i(t)|^2||F_i||_{L^2(0,1)}^2+\frac{1}{2}||W_{i,m}(t)||_{L^2(0,1)}^2.$$
    Moreover, 
    $$\left|\int_0^1\frac{\partial b}{\partial \xi_i}|W_{i,m}|^2d\xi_i\right|\leq \left\|\frac{\partial b}{\partial \xi_i}(\cdot,t)\right\|_{L^{\infty}(0,1)}||W_{i,m}||_{L^2(0,1)}^2.$$
    Substitution all these estimations into \eqref{eq63} gives the energy identity
    $$\frac{d}{dt}||W_{i,m}(t)||_{L^2(0,1)}^2+2a_i(t)\left\|\frac{\partial W_{i,m}}{\partial\xi_i}(t)\right\|_{L^2(0,1)}^2\leq\left(1+\left\|\frac{\partial b}{\partial \xi_i}(\cdot,t)\right\|_{L^{\infty}(0,1)}\right)||W_{i,m}(t)||_{L^2(0,1)}^2$$
    $$+|\rho_i(t)|||F_i(t)||_{L^2(0,1)}^2.$$
    As it is shown that $\rho_i(t)\equiv 0$, by integrating over $[0,t)$ with $t\in [0,T)$, we can obtain the next estimation
    $$||W_{i,m}(t)||_{L^2(0,1)}^2+2N_i\int_0^T\left\|\frac{\partial W_{i,m}}{\partial\xi_i}(\tau)\right\|_{L^2(0,1)}^2d\tau \leq C_0\int_0^T ||W_{i,m}(t)||_{L^2(0,1)}^2$$,
    where $C_0=\left|1+\left\|\frac{\partial b}{\partial \xi_i}(\cdot,t)\right\|_{L^{\infty}(0,1)}\right|$. Hence, by Gronwall's inequality we can conclude that $||W_{i,m}(t)||_{L^2(0,1)}^2\equiv 0$, which leads to $W_{i,m}=0$ and $V_{i,m}^{(1)}=V_{i,m}^{(2)}$.
\end{proof}

\end{document}
